\documentclass[11pt,a4paper]{article}

% Core packages
\usepackage[margin=1in]{geometry}
\usepackage[utf8]{inputenc}
\usepackage[T1]{fontenc}
\usepackage{graphicx}
\usepackage[colorlinks=true, linkcolor=blue, urlcolor=blue, citecolor=blue]{hyperref}
\usepackage[table]{xcolor}
\usepackage{colortbl}
\usepackage{booktabs}
\usepackage{tabularx}
\usepackage{float}
\usepackage{enumitem}
\usepackage{longtable}
\usepackage{amsmath, amssymb}
\usepackage[most]{tcolorbox}
\usepackage{titlesec}
\usepackage{fancyhdr}
\usepackage{microtype}
\usepackage{parskip}

% Widow/orphan control
\widowpenalty=10000
\clubpenalty=10000

% Global list compaction
\setlist[itemize]{nosep, leftmargin=*, topsep=2pt, partopsep=0pt}
\setlist[enumerate]{nosep, leftmargin=*, topsep=2pt, partopsep=0pt}

% Section heading style
\titleformat{\section}{\large\bfseries\color{blue!70!black}}{}{0em}{\thesection\quad}[\titlerule]
\titleformat{\subsection}{\normalsize\bfseries}{}{0em}{\thesubsection\quad}

% Header / footer
\pagestyle{fancy}
\fancyhf{}
\lhead{\small DAP391m --- Group 8}
\rhead{\small Updated Project Pipeline}
\cfoot{\thepage}
\renewcommand{\headrulewidth}{0.4pt}

% tcolorbox styles
\tcbset{
  warningbox/.style={
    colback=red!5, colframe=red!65!black,
    fonttitle=\bfseries, title={Critical Warning},
    left=6pt, right=6pt, top=4pt, bottom=4pt
  },
  infobox/.style={
    colback=blue!5, colframe=blue!60!black,
    fonttitle=\bfseries,
    left=6pt, right=6pt, top=4pt, bottom=4pt
  },
  adoptbox/.style={
    colback=green!5, colframe=green!55!black,
    fonttitle=\bfseries,
    left=6pt, right=6pt, top=4pt, bottom=4pt
  },
  decisionbox/.style={
    colback=orange!5, colframe=orange!70!black,
    fonttitle=\bfseries,
    left=6pt, right=6pt, top=4pt, bottom=4pt
  }
}

% Column type for wide tables
\newcolumntype{L}[1]{>{\raggedright\arraybackslash}p{#1}}

\begin{document}

% ──────────────────────────────────────────
%  TITLE PAGE
% ──────────────────────────────────────────
\begin{titlepage}
\centering
\vspace*{2cm}
{\LARGE\bfseries Updated Project Pipeline\par}
\vspace{0.5em}
{\large\color{blue!70!black} Supplier Lead-Time Risk Prediction --- DAP391m, Summer 2026\par}
\vspace{0.3em}
{\normalsize Group 8\par}
\vspace{2cm}

\begin{tabular}{ccc}
  Nguyen Hoai Khanh & Ho Lam Bao Dang & Duong Gia Bao \\
  \small K@itsdre.me & \small baodang10506@gmail.com & \small ilyrabao@gmail.com
\end{tabular}

\vspace{0.8cm}
{\small FPT University HCMC \\ Supervisor: Mr.\ Nguyen Hoai Linh\par}
\vspace{1.5cm}

\begin{tcolorbox}[infobox, title={Document Purpose}]
This document records every pipeline decision made after receiving supervisor feedback
(``data not good enough'') and reviewing the reference paper by Orajaka \& Okolie (2025).
It specifies exactly what to adopt, what to keep, what to drop, and how each pipeline
script must change. No code is modified by this document; it is the specification that
guides code changes.
\end{tcolorbox}

\vfill
{\small Version 1.0 --- 2026-05-20\par}
\end{titlepage}

\tableofcontents
\newpage

% ──────────────────────────────────────────
\section{Executive Summary}
% ──────────────────────────────────────────

The supervisor's assessment that the original data was ``ok but not good enough'' was
confirmed by quantitative analysis: the original joined dataset produced only 704 rows
from 7 unique suppliers, providing insufficient coverage for reliable generalization.
The team sourced a replacement dataset (\texttt{supply\_chain\_risk\_dataset.csv})
containing 2,478 rows across 51 suppliers, which directly resolves the scale concern.

In parallel, the reference paper---\textit{Using Artificial Intelligence to predict and
optimize supplier lead times in procurement operations} (Orajaka \& Okolie, 2025)---was
reviewed in full. While the paper adopts a regression framing that ultimately fails
($R^2 = -0.09$), it contributes three high-value techniques: expanding-window supplier
history features, temporal order features, and a rigorous theoretical framework spanning
five management theories.

The central update is therefore two-fold: migrate to the new dataset and integrate the
paper's best feature-engineering practices. The project's classification framing,
four-model benchmark, PR-AUC primary metric, and stratified cross-validation are all
retained---they are demonstrably stronger than the paper's approach and better suited
to the procurement risk-alerting use case.

% ──────────────────────────────────────────
\section{Dataset Transition}
% ──────────────────────────────────────────

\subsection{Quantitative Comparison}

Table~\ref{tab:dataset-compare} presents a side-by-side comparison of the original
three-CSV dataset and the new single-CSV dataset across the dimensions most relevant
to model reliability.

\begin{table}[htbp]
\centering
\caption{Original vs.\ new dataset --- key metrics}
\label{tab:dataset-compare}
\renewcommand{\arraystretch}{1.25}
\begin{tabularx}{\textwidth}{L{4.2cm} X X}
\toprule
\textbf{Dimension} & \textbf{Original (3 CSVs)} & \textbf{New Dataset} \\
\midrule
\rowcolor{gray!8}
Rows after join & 704 & 2,478 \\
Unique suppliers & 7 & 51 \\
\rowcolor{gray!8}
Missing values & Present & Zero \\
Target variable & \texttt{delivery\_status} (binary) & \texttt{risk\_label} (3-class) \\
\rowcolor{gray!8}
Class distribution & 16\% Delayed / 84\% On-Time & 63\% High / 19\% Low / 18\% Medium \\
Supplier features & \texttt{lead\_time\_days}, \texttt{satisfaction\_score} & \texttt{supplier\_lead\_time\_days}, \texttt{supplier\_quality\_score}, \texttt{supplier\_reliability\_index} \\
\rowcolor{gray!8}
Environmental features & None & Port delay, weather disruption, fuel price, market demand \\
Machine/IoT features & None & Temperature, vibration, runtime \\
\rowcolor{gray!8}
Temporal coverage & Multiple years & Mar--Jul 2025 (5 months) \\
\bottomrule
\end{tabularx}
\end{table}

\subsection{Target Variable Decision}

The new dataset provides two target columns: \texttt{risk\_probability} (continuous
0--1) and \texttt{risk\_label} (categorical: High/Medium/Low). A correlation analysis
revealed that \texttt{risk\_probability} correlates at $r = 0.898$ with
\texttt{port\_delay\_days}, confirming the label is synthetically derived rather than
observed from real delivery outcomes. The detailed implications are addressed in
Section~\ref{sec:leakage}.

For modelling, the project will use \texttt{risk\_label} as the primary target in its
three-class form (High, Medium, Low). A binary collapse (High vs.\ not-High) will be
computed as a secondary comparison target to maintain continuity with the original
proposal's binary framing.

% ──────────────────────────────────────────
\section{Data Quality Warning: Leakage Risk}
\label{sec:leakage}
% ──────────────────────────────────────────

\begin{tcolorbox}[warningbox]
\textbf{Finding:} \texttt{port\_delay\_days} correlates with \texttt{risk\_probability}
at $r = 0.898$. All other features correlate at $r \leq 0.102$. The target label is
therefore an algebraic function of \texttt{port\_delay\_days}, not an independent
business outcome.
\end{tcolorbox}

This finding has two direct pipeline consequences. First, \texttt{port\_delay\_days}
must be excluded from the model feature set entirely. A model trained with it as a
feature would learn a trivial mapping and achieve near-perfect accuracy on a proxy
signal rather than a real procurement risk. Second, \texttt{risk\_probability} must
also be excluded for the same reason---it \textit{is} the target transformed.

The remaining 12 operational and environmental features (temperature, vibration,
runtime, inventory level, pending orders, supplier lead time, supplier quality score,
supplier reliability index, fuel price index, market demand index, and weather
disruption score) represent genuine supply chain signals. After adding the engineered
features described in Section~\ref{sec:feature-eng}, the final feature matrix will
contain approximately 17--20 columns.

The synthetic nature of the target must be disclosed in the project report's data
section. A suggested sentence: \textit{``The \texttt{risk\_label} target variable is
derived primarily from port delay conditions ($r = 0.898$), making it a proxy risk
score rather than a directly observed delivery outcome. To prevent target leakage,
\texttt{port\_delay\_days} is excluded from the feature set.''}

% ──────────────────────────────────────────
\section{Updated Feature Engineering Plan}
\label{sec:feature-eng}
% ──────────────────────────────────────────

\subsection{Features Adopted from the Paper}

\textbf{Expanding-Window Supplier History.} The paper's second and third most important
SHAP features are \texttt{supplier\_avg\_lead} and \texttt{supplier\_std\_lead},
computed using an expanding window per supplier. These capture a supplier's historical
reliability: the mean quantifies average performance, while the standard deviation
captures consistency. Both will be engineered from \texttt{supplier\_lead\_time\_days}
in the new dataset. A critical implementation detail is the \texttt{.shift(1)} offset,
which ensures the current row's lead-time value does not contribute to its own
prediction feature, thereby preventing look-ahead leakage.

\begin{tcolorbox}[adoptbox, title={Expanding Window --- Implementation Spec}]
\begin{verbatim}
df = df.sort_values('timestamp')
df['supplier_avg_lead'] = df.groupby('supplier_id')['supplier_lead_time_days'] \
    .transform(lambda x: x.expanding().mean().shift(1))
df['supplier_std_lead'] = df.groupby('supplier_id')['supplier_lead_time_days'] \
    .transform(lambda x: x.expanding().std().shift(1))
\end{verbatim}
With $\approx$48.6 observations per supplier on average, the expanding windows will
have sufficient history to stabilise by the 10th observation per supplier.
\end{tcolorbox}

\textbf{Temporal Features.} The paper extracts \texttt{order\_month} and
\texttt{order\_weekday} from the order date to capture seasonal and scheduling patterns.
The new dataset's \texttt{timestamp} column enables equivalent extraction plus an
\texttt{order\_hour} feature that reflects shift-level operational patterns---a
dimension the paper did not exploit.

\subsection{Original Project Features to Engineer}

\textbf{Supplier Risk Score.} Combining multiple supplier signals into a single
composite score makes the signal denser and provides an interpretable KPI for the
Power BI dashboard. The composite weights supplier lead-time history (40\%), lead-time
consistency (30\%), quality score (20\%), and reliability index (10\%). These weights
are heuristic starting points; SHAP analysis after modelling will validate whether
the composite or its constituent parts better explain predictions.

\begin{tcolorbox}[infobox, title={Supplier Risk Score --- Implementation Spec}]
\begin{verbatim}
df['supplier_risk_score'] = (
    df['supplier_avg_lead'] * 0.40 +
    df['supplier_std_lead'] * 0.30 +
    (100 - df['supplier_quality_score']) / 100 * 0.20 +
    (1 - df['supplier_reliability_index']) * 0.10
)
\end{verbatim}
Rows where \texttt{supplier\_avg\_lead} is NaN (first observation per supplier) should
have \texttt{supplier\_risk\_score} set to NaN and subsequently imputed with the
supplier-level median.
\end{tcolorbox}

\textbf{External Risk Composite.} The paper explicitly identifies environmental factors
as a gap in the literature. The new dataset fills this gap with weather disruption,
fuel price, and market demand. These will be combined into a single external disruption
index (after \texttt{port\_delay\_days} is excluded). This composite provides an
interpretable environmental-pressure signal for the Streamlit app.

\subsection{Encoding Changes}

The original pipeline used \texttt{LabelEncoder} for \texttt{supplier\_id}. This will
be replaced with One-Hot Encoding for \texttt{supplier\_id} (51 unique values = 50
binary columns, manageable). \texttt{machine\_id} has 201 unique values; OHE would
produce 200 noisy indicator columns. Instead, \texttt{machine\_id} will be dropped
entirely---machine-level identity is not a procurement variable and the supplier
features already capture the relevant supplier-level variation.

\subsection{Feature Inventory: Final Matrix}

Table~\ref{tab:features} summarises the full feature set after engineering and encoding.

\begin{table}[htbp]
\centering
\caption{Final feature matrix --- new dataset after engineering}
\label{tab:features}
\renewcommand{\arraystretch}{1.2}
\begin{tabularx}{\textwidth}{L{4.5cm} L{2.5cm} X}
\toprule
\textbf{Feature} & \textbf{Source} & \textbf{Notes} \\
\midrule
\rowcolor{gray!8}
\texttt{temperature\_C} & Raw & Machine operational signal \\
\texttt{vibration\_level} & Raw & Machine operational signal \\
\rowcolor{gray!8}
\texttt{machine\_runtime\_hours} & Raw & Machine utilisation \\
\texttt{inventory\_level\_units} & Raw & Inventory pressure \\
\rowcolor{gray!8}
\texttt{pending\_orders} & Raw & Demand pressure \\
\texttt{supplier\_lead\_time\_days} & Raw & Core supplier metric \\
\rowcolor{gray!8}
\texttt{supplier\_quality\_score} & Raw & Proxy for defect rate \\
\texttt{supplier\_reliability\_index} & Raw & Delivery consistency \\
\rowcolor{gray!8}
\texttt{fuel\_price\_index} & Raw & External cost pressure \\
\texttt{weather\_disruption\_score} & Raw & Environmental disruption \\
\rowcolor{gray!8}
\texttt{market\_demand\_index} & Raw & Demand-side pressure \\
\texttt{order\_month} & Engineered & Seasonal pattern \\
\rowcolor{gray!8}
\texttt{order\_weekday} & Engineered & Scheduling pattern \\
\texttt{order\_hour} & Engineered & Shift-level pattern \\
\rowcolor{gray!8}
\texttt{supplier\_avg\_lead} & Engineered & Historical mean (paper) \\
\texttt{supplier\_std\_lead} & Engineered & Historical variability (paper) \\
\rowcolor{gray!8}
\texttt{supplier\_risk\_score} & Engineered & Composite (new) \\
\texttt{external\_risk\_score} & Engineered & Environmental composite \\
\rowcolor{gray!8}
\texttt{supplier\_id\_*} (OHE) & Encoded & 50 binary columns \\
\midrule
\texttt{port\_delay\_days} & \textcolor{red!70!black}{\textbf{DROPPED}} & Near-target leakage ($r = 0.898$) \\
\texttt{risk\_probability} & \textcolor{red!70!black}{\textbf{DROPPED}} & Direct target leakage \\
\texttt{machine\_id} & \textcolor{red!70!black}{\textbf{DROPPED}} & 201 categories, no procurement signal \\
\texttt{timestamp} & \textcolor{red!70!black}{\textbf{DROPPED}} & Replaced by temporal features \\
\bottomrule
\end{tabularx}
\end{table}

% ──────────────────────────────────────────
\section{Pipeline Script Specifications}
% ──────────────────────────────────────────

Each of the six pipeline scripts requires updates for the new dataset. This section
specifies exactly what changes in each file. No code is modified here; these are
specifications for the developer to implement.

\subsection{Script 01 --- Ingestion and Cleaning}

The script must be updated to load \texttt{Data/supply\_chain\_risk\_dataset.csv}
instead of the three original CSVs. The three-way join logic (customer $\rightarrow$
shipment $\rightarrow$ logistics) is replaced by single-file ingestion. Type coercion
should parse \texttt{timestamp} as \texttt{datetime64}, and \texttt{risk\_label} as a
categorical with ordered levels (Low $<$ Medium $<$ High). The null check will pass
trivially since the new dataset has zero missing values, but the check should remain
as a pipeline guard. Output: \texttt{Data/filtered/clean\_data.csv}.

\subsection{Script 02 --- SQL Analysis}

The SQL queries must be rewritten for the new schema. The six required queries map to
the new columns as follows: average lead time per supplier uses
\texttt{supplier\_lead\_time\_days} grouped by \texttt{supplier\_id}; delay frequency
uses \texttt{risk\_label = 'High'} as the delay proxy; volume--delay correlation
compares \texttt{pending\_orders} against \texttt{risk\_label}; monthly trends group by
\texttt{strftime('\%Y-\%m', timestamp)}; year-over-year growth uses the same timestamp
grouping at annual resolution; and the penetration index computes the fraction of
High-risk events per supplier relative to the dataset mean.

\subsection{Script 03 --- EDA}

The EDA script should add three new analyses: a correlation heatmap confirming the
\texttt{port\_delay\_days} leakage signal (to document it explicitly), a class
balance chart for the three-class target, and a per-supplier box plot of
\texttt{supplier\_lead\_time\_days} to visualise the 51-supplier spread (analogous to
Figure~2 in the reference paper).

\subsection{Script 04 --- Feature Engineering}

This is the most substantially changed script. The changes are:
\begin{enumerate}
  \item Parse timestamp and extract \texttt{order\_month}, \texttt{order\_weekday}, \texttt{order\_hour}.
  \item Drop \texttt{port\_delay\_days}, \texttt{risk\_probability}, \texttt{machine\_id}, \texttt{timestamp}.
  \item Sort by \texttt{timestamp} and compute expanding-window supplier features with shift.
  \item Compute \texttt{supplier\_risk\_score} and \texttt{external\_risk\_score}.
  \item Apply IQR capping to continuous features (retain existing IQR logic).
  \item Replace \texttt{LabelEncoder} with \texttt{pd.get\_dummies} for \texttt{supplier\_id}.
  \item Encode \texttt{risk\_label} as ordered integer (Low=0, Medium=1, High=2) for multi-class models.
  \item Stratified train/test split (80/20) preserving class proportions.
\end{enumerate}

\subsection{Script 05 --- Modelling}

The four classifiers (Logistic Regression, Decision Tree, Random Forest, XGBoost) are
retained without change to their architecture. The target changes from binary to
three-class, which requires two configuration updates: \texttt{multi\_class='ovr'} for
Logistic Regression and \texttt{objective='multi:softprob'} with
\texttt{num\_class=3} for XGBoost. The stratified five-fold cross-validation
configuration is unchanged. Primary evaluation metrics shift from binary PR-AUC to
macro-averaged PR-AUC across the three classes, with per-class recall reported
separately. A binary secondary analysis (High vs.\ not-High) should be run alongside
the three-class models to maintain comparability with the original proposal.

\subsection{Script 06 --- Advanced Visualisation}

The Folium risk heatmap should be adapted to use \texttt{supplier\_id} as the grouping
key. Since the new dataset lacks geographic coordinates, a simulated spatial layout
(supplier IDs mapped to a grid) or a bubble chart by risk level provides an equivalent
visual. The existing Plotly boxplots by supplier and lead-time histograms translate
directly to the new schema.

% ──────────────────────────────────────────
\section{Model Evaluation Framework}
% ──────────────────────────────────────────

\subsection{Primary Metrics}

Macro-averaged PR-AUC is the primary metric for the three-class setting, ensuring that
performance on the minority classes (Medium and Low) is not obscured by the dominant
High class. Per-class recall is reported separately so that the High-risk recall---the
most operationally critical figure---can be evaluated in isolation. This mirrors the
original proposal's emphasis on recall for the delayed class.

\subsection{Secondary Metrics}

ROC-AUC (macro OvR), weighted F1-score, and the three-class confusion matrix serve as
secondary metrics. For the binary secondary analysis, the full set of original secondary
metrics (binary PR-AUC, binary ROC-AUC, F1, confusion matrix) is preserved.

\subsection{Validation}

Stratified five-fold cross-validation is retained. With 2,478 rows, each fold contains
approximately 496 samples, which provides stable variance estimates across all three
classes. The final test set (held-out 20\%) is used only for final reporting; all
hyperparameter decisions are made on the cross-validation folds.

\begin{tcolorbox}[decisionbox, title={Why Not Switch to Regression?}]
The reference paper frames the task as regression (predicting exact lead-time days)
and achieves $R^2 = -0.09$, meaning the model performs worse than a simple mean
predictor. Procurement teams need an actionable risk flag (High / Medium / Low), not a
day-count estimate. The classification framing is retained for both practical and
performance reasons.
\end{tcolorbox}

% ──────────────────────────────────────────
\section{Theoretical Framework for the Report}
% ──────────────────────────────────────────

The reference paper grounds its work in five management theories. All five directly
justify the AI-driven classification approach used in this project and should be cited
in the final report's literature review section. Table~\ref{tab:theories} maps each
theory to its specific justification for this project.

\begin{table}[htbp]
\centering
\caption{Theoretical framework adopted from Orajaka \& Okolie (2025)}
\label{tab:theories}
\renewcommand{\arraystretch}{1.3}
\begin{tabularx}{\textwidth}{L{4cm} X}
\toprule
\textbf{Theory} & \textbf{Justification for this Project} \\
\midrule
\rowcolor{gray!8}
Transaction Cost Economics (Williamson, 1979) &
  Lead-time delays are hidden transaction costs. ML quantifies this cost and enables
  procurement teams to price supplier risk into sourcing decisions. \\
Resource-Based View (Barney, 1991) &
  Supplier reliability is a strategic, firm-specific resource. Predictive modelling
  identifies which suppliers constitute a competitive advantage and which are a
  liability. \\
\rowcolor{gray!8}
Lean Supply Chain Theory (Womack \& Jones, 1996) &
  Delay prediction enables proactive waste elimination. By anticipating late deliveries,
  procurement can adjust safety-stock policies and avoid production stoppages. \\
Complex Adaptive Systems (Choi et al., 2001) &
  Supply chains are non-linear, emergent networks. Tree-based ensemble models (RF,
  XGBoost) are better suited to this non-linearity than linear regression. \\
\rowcolor{gray!8}
Decision Theory + Risk Management (Raiffa, 1968) &
  Probabilistic risk scores from the classifier support rational procurement
  decisions under uncertainty, enabling cost--risk trade-off analysis. \\
\bottomrule
\end{tabularx}
\end{table}

% ──────────────────────────────────────────
\section{Full Adoption Decision Table}
% ──────────────────────────────────────────

Table~\ref{tab:adoption} is the master decision record for this pipeline update,
covering every element evaluated across the original dataset, the new dataset, and
the reference paper.

\begin{longtable}{L{5.2cm} L{2.2cm} L{5.8cm}}
\caption{Master adoption table --- pipeline update decisions}
\label{tab:adoption} \\
\toprule
\textbf{Element} & \textbf{Action} & \textbf{Rationale} \\
\midrule
\endfirsthead
\multicolumn{3}{l}{\small\textit{Table~\ref{tab:adoption} continued}} \\
\toprule
\textbf{Element} & \textbf{Action} & \textbf{Rationale} \\
\midrule
\endhead
\midrule
\multicolumn{3}{r}{\small\textit{Continued on next page}} \\
\endfoot
\bottomrule
\endlastfoot
\rowcolor{green!8}
\texttt{supplier\_avg\_lead} via expanding window & \textbf{ADOPT} & Paper's 2nd SHAP feature; captures supplier history \\
\texttt{supplier\_std\_lead} via expanding window & \textbf{ADOPT} & Paper's 3rd SHAP feature; captures delivery consistency \\
\rowcolor{green!8}
One-Hot Encoding for \texttt{supplier\_id} & \textbf{ADOPT} & Avoids false ordinal relationships in tree models \\
\texttt{order\_month}, \texttt{order\_weekday}, \texttt{order\_hour} & \textbf{ADOPT} & Temporal features from paper; new dataset enables \texttt{order\_hour} \\
\rowcolor{green!8}
Five theoretical frameworks in literature review & \textbf{ADOPT} & Directly justifies ML approach; paper provides full citations \\
IQR outlier capping on continuous features & \textbf{RETAIN (extend)} & Already in pipeline; extend to new dataset columns \\
\rowcolor{blue!5}
Classification framing (not regression) & \textbf{KEEP} & Paper's regression failed ($R^2 = -0.09$); classification is more actionable \\
Four-model benchmarking & \textbf{KEEP} & More rigorous than paper's single-model approach \\
\rowcolor{blue!5}
PR-AUC + Recall as primary metrics & \textbf{KEEP} & Correct for imbalanced multi-class; paper uses inappropriate MAE/RMSE \\
Stratified 5-fold cross-validation & \textbf{KEEP} & Paper has no CV; provides variance estimates \\
\rowcolor{blue!5}
Per-prediction SHAP waterfall in Streamlit & \textbf{KEEP} & Paper shows aggregate SHAP only; per-prediction is the differentiator \\
Power BI supplier scorecard & \textbf{KEEP} & Course requirement; not present in paper \\
\rowcolor{blue!5}
SQL 6-query analysis & \textbf{KEEP} & Course requirement; queries rewritten for new schema \\
3-class target (High/Medium/Low) & \textbf{ADD} & More nuanced than binary; aligns with Power BI risk tiers \\
\rowcolor{yellow!15}
\texttt{supplier\_risk\_score} composite & \textbf{ADD} & Paper-inspired; provides interpretable KPI for dashboard \\
\texttt{external\_risk\_score} composite & \textbf{ADD} & Addresses paper's stated future-work gap \\
\rowcolor{yellow!15}
Binary secondary analysis (High vs.\ not-High) & \textbf{ADD} & Maintains comparability with original proposal \\
Data leakage disclosure in report & \textbf{ADD} & $r = 0.898$ finding must be documented \\
\rowcolor{red!8}
\texttt{port\_delay\_days} as model feature & \textbf{DROP} & Near-target leakage ($r = 0.898$) \\
\texttt{risk\_probability} as model feature & \textbf{DROP} & Direct target leakage \\
\rowcolor{red!8}
\texttt{machine\_id} OHE & \textbf{DROP} & 201 categories; no procurement signal \\
Regression task (MAE / RMSE / $R^2$) & \textbf{REJECT} & Paper demonstrated failure; wrong framing \\
\rowcolor{red!8}
Original 3-CSV join pipeline & \textbf{RETIRE} & Replaced by new dataset; 7 suppliers insufficient \\
\end{longtable}

% ──────────────────────────────────────────
\section{Deliverables Status}
% ──────────────────────────────────────────

\subsection{Course Requirements}

All course deliverables are unaffected by the dataset change. The four-model
comparison, SQL analysis file (\texttt{sql/analysis.sql}), AI Audit Log
(\texttt{docs/AI\_AuditLog\_Template\_DAP391m.xlsx}), and final LaTeX report
(\texttt{report/main.tex}) remain as specified. The modelling target changes from
binary to three-class, which updates the classifier configurations and evaluation
tables but does not change the deliverable categories.

\subsection{Applied Deliverables}

The Power BI supplier scorecard gains a natural upgrade path: the three-class risk
label (High/Medium/Low) maps directly to traffic-light risk alerts without requiring
threshold tuning. The Streamlit application's SHAP waterfall explanation is unchanged;
the input fields will be updated to match the new feature set.

\subsection{Report Literature Review}

A minimum of three citations from Orajaka \& Okolie (2025) should be incorporated into
the final report: the feature importance finding that supplier reliability dominates
lead-time prediction, the gap identification for environmental factors (now addressed
by the new dataset), and the recommendation for ensemble over linear models. These
citations strengthen the report's positioning relative to prior work.

% ──────────────────────────────────────────
\section{Implementation Order}
% ──────────────────────────────────────────

The pipeline updates should be implemented in the following sequence to minimise
integration risk. Scripts 01 and 04 carry the highest change surface; scripts 05 and
06 require only configuration updates once the feature matrix is validated.

\begin{enumerate}
  \item \textbf{Script 01} --- Update ingestion for new CSV; validate shape (2,478 $\times$ 17) and zero nulls.
  \item \textbf{Script 03} --- Run EDA on raw new data; confirm leakage correlation visually.
  \item \textbf{Script 04} --- Implement all feature engineering changes; validate final matrix shape.
  \item \textbf{Script 02} --- Rewrite SQL queries for new schema; validate output tables.
  \item \textbf{Script 05} --- Update classifier configs for three-class; run CV; report metrics.
  \item \textbf{Script 06} --- Update visualisations for new schema; generate report figures.
  \item \textbf{Streamlit app} --- Update input fields and SHAP waterfall for new feature set.
  \item \textbf{Power BI} --- Rebuild scorecard with three-class risk tiers.
\end{enumerate}

Running the pipeline in this order ensures that each downstream script receives a
validated input before it is updated, making integration errors traceable to a single
script at a time.

\end{document}
